\IfFileExists{stacks-project.cls}{%
\documentclass{stacks-project}
}{%
\documentclass{amsart}
}

% For dealing with references we use the comment environment
\usepackage{verbatim}
\newenvironment{reference}{\comment}{\endcomment}
%\newenvironment{reference}{}{}
\newenvironment{slogan}{\comment}{\endcomment}
\newenvironment{history}{\comment}{\endcomment}

% For commutative diagrams we use Xy-pic
\usepackage[all]{xy}

% We use 2cell for 2-commutative diagrams.
\xyoption{2cell}
\UseAllTwocells

% We use multicol for the list of chapters between chapters
\usepackage{multicol}

% This is generally recommended for better output
\usepackage{lmodern}
\usepackage[T1]{fontenc}

% For cross-file-references
\usepackage{xr-hyper}

% Package for hypertext links:
\usepackage{hyperref}

% For any local file, say "hello.tex" you want to link to please
% use \externaldocument[hello-]{hello}
\externaldocument[introduction-]{introduction}
\externaldocument[conventions-]{conventions}
\externaldocument[sets-]{sets}
\externaldocument[categories-]{categories}
\externaldocument[topology-]{topology}
\externaldocument[sheaves-]{sheaves}
\externaldocument[sites-]{sites}
\externaldocument[stacks-]{stacks}
\externaldocument[fields-]{fields}
\externaldocument[algebra-]{algebra}
\externaldocument[brauer-]{brauer}
\externaldocument[homology-]{homology}
\externaldocument[derived-]{derived}
\externaldocument[simplicial-]{simplicial}
\externaldocument[more-algebra-]{more-algebra}
\externaldocument[smoothing-]{smoothing}
\externaldocument[modules-]{modules}
\externaldocument[sites-modules-]{sites-modules}
\externaldocument[injectives-]{injectives}
\externaldocument[cohomology-]{cohomology}
\externaldocument[sites-cohomology-]{sites-cohomology}
\externaldocument[dga-]{dga}
\externaldocument[dpa-]{dpa}
\externaldocument[sdga-]{sdga}
\externaldocument[hypercovering-]{hypercovering}
\externaldocument[schemes-]{schemes}
\externaldocument[constructions-]{constructions}
\externaldocument[properties-]{properties}
\externaldocument[morphisms-]{morphisms}
\externaldocument[coherent-]{coherent}
\externaldocument[divisors-]{divisors}
\externaldocument[limits-]{limits}
\externaldocument[varieties-]{varieties}
\externaldocument[topologies-]{topologies}
\externaldocument[descent-]{descent}
\externaldocument[perfect-]{perfect}
\externaldocument[more-morphisms-]{more-morphisms}
\externaldocument[flat-]{flat}
\externaldocument[groupoids-]{groupoids}
\externaldocument[more-groupoids-]{more-groupoids}
\externaldocument[etale-]{etale}
\externaldocument[chow-]{chow}
\externaldocument[intersection-]{intersection}
\externaldocument[pic-]{pic}
\externaldocument[weil-]{weil}
\externaldocument[adequate-]{adequate}
\externaldocument[dualizing-]{dualizing}
\externaldocument[duality-]{duality}
\externaldocument[discriminant-]{discriminant}
\externaldocument[derham-]{derham}
\externaldocument[local-cohomology-]{local-cohomology}
\externaldocument[algebraization-]{algebraization}
\externaldocument[curves-]{curves}
\externaldocument[resolve-]{resolve}
\externaldocument[models-]{models}
\externaldocument[functors-]{functors}
\externaldocument[equiv-]{equiv}
\externaldocument[pione-]{pione}
\externaldocument[etale-cohomology-]{etale-cohomology}
\externaldocument[proetale-]{proetale}
\externaldocument[relative-cycles-]{relative-cycles}
\externaldocument[more-etale-]{more-etale}
\externaldocument[trace-]{trace}
\externaldocument[crystalline-]{crystalline}
\externaldocument[spaces-]{spaces}
\externaldocument[spaces-properties-]{spaces-properties}
\externaldocument[spaces-morphisms-]{spaces-morphisms}
\externaldocument[decent-spaces-]{decent-spaces}
\externaldocument[spaces-cohomology-]{spaces-cohomology}
\externaldocument[spaces-limits-]{spaces-limits}
\externaldocument[spaces-divisors-]{spaces-divisors}
\externaldocument[spaces-over-fields-]{spaces-over-fields}
\externaldocument[spaces-topologies-]{spaces-topologies}
\externaldocument[spaces-descent-]{spaces-descent}
\externaldocument[spaces-perfect-]{spaces-perfect}
\externaldocument[spaces-more-morphisms-]{spaces-more-morphisms}
\externaldocument[spaces-flat-]{spaces-flat}
\externaldocument[spaces-groupoids-]{spaces-groupoids}
\externaldocument[spaces-more-groupoids-]{spaces-more-groupoids}
\externaldocument[bootstrap-]{bootstrap}
\externaldocument[spaces-pushouts-]{spaces-pushouts}
\externaldocument[spaces-chow-]{spaces-chow}
\externaldocument[groupoids-quotients-]{groupoids-quotients}
\externaldocument[spaces-more-cohomology-]{spaces-more-cohomology}
\externaldocument[spaces-simplicial-]{spaces-simplicial}
\externaldocument[spaces-duality-]{spaces-duality}
\externaldocument[formal-spaces-]{formal-spaces}
\externaldocument[restricted-]{restricted}
\externaldocument[spaces-resolve-]{spaces-resolve}
\externaldocument[formal-defos-]{formal-defos}
\externaldocument[defos-]{defos}
\externaldocument[cotangent-]{cotangent}
\externaldocument[examples-defos-]{examples-defos}
\externaldocument[algebraic-]{algebraic}
\externaldocument[examples-stacks-]{examples-stacks}
\externaldocument[stacks-sheaves-]{stacks-sheaves}
\externaldocument[criteria-]{criteria}
\externaldocument[artin-]{artin}
\externaldocument[quot-]{quot}
\externaldocument[stacks-properties-]{stacks-properties}
\externaldocument[stacks-morphisms-]{stacks-morphisms}
\externaldocument[stacks-limits-]{stacks-limits}
\externaldocument[stacks-cohomology-]{stacks-cohomology}
\externaldocument[stacks-perfect-]{stacks-perfect}
\externaldocument[stacks-introduction-]{stacks-introduction}
\externaldocument[stacks-more-morphisms-]{stacks-more-morphisms}
\externaldocument[stacks-geometry-]{stacks-geometry}
\externaldocument[moduli-]{moduli}
\externaldocument[moduli-curves-]{moduli-curves}
\externaldocument[examples-]{examples}
\externaldocument[exercises-]{exercises}
\externaldocument[guide-]{guide}
\externaldocument[desirables-]{desirables}
\externaldocument[coding-]{coding}
\externaldocument[obsolete-]{obsolete}
\externaldocument[fdl-]{fdl}
\externaldocument[index-]{index}

% Theorem environments.
%
\theoremstyle{plain}
\newtheorem{theorem}[subsection]{Theorem}
\newtheorem{proposition}[subsection]{Proposition}
\newtheorem{lemma}[subsection]{Lemma}

\theoremstyle{definition}
\newtheorem{definition}[subsection]{Definition}
\newtheorem{example}[subsection]{Example}
\newtheorem{exercise}[subsection]{Exercise}
\newtheorem{situation}[subsection]{Situation}

\theoremstyle{remark}
\newtheorem{remark}[subsection]{Remark}
\newtheorem{remarks}[subsection]{Remarks}

\numberwithin{equation}{subsection}

% Macros
%
\def\lim{\mathop{\mathrm{lim}}\nolimits}
\def\colim{\mathop{\mathrm{colim}}\nolimits}
\def\Spec{\mathop{\mathrm{Spec}}}
\def\Hom{\mathop{\mathrm{Hom}}\nolimits}
\def\Ext{\mathop{\mathrm{Ext}}\nolimits}
\def\SheafHom{\mathop{\mathcal{H}\!\mathit{om}}\nolimits}
\def\SheafExt{\mathop{\mathcal{E}\!\mathit{xt}}\nolimits}
\def\Sch{\mathit{Sch}}
\def\Mor{\mathop{\mathrm{Mor}}\nolimits}
\def\Ob{\mathop{\mathrm{Ob}}\nolimits}
\def\Sh{\mathop{\mathit{Sh}}\nolimits}
\def\NL{\mathop{N\!L}\nolimits}
\def\CH{\mathop{\mathrm{CH}}\nolimits}
\def\proetale{{pro\text{-}\acute{e}tale}}
\def\etale{{\acute{e}tale}}
\def\QCoh{\mathit{QCoh}}
\def\Ker{\mathop{\mathrm{Ker}}}
\def\Im{\mathop{\mathrm{Im}}}
\def\Coker{\mathop{\mathrm{Coker}}}
\def\Coim{\mathop{\mathrm{Coim}}}

% Boxtimes
%
\DeclareMathSymbol{\boxtimes}{\mathbin}{AMSa}{"02}

%
% Macros for moduli stacks/spaces
%
\def\QCohstack{\mathcal{QC}\!\mathit{oh}}
\def\Cohstack{\mathcal{C}\!\mathit{oh}}
\def\Spacesstack{\mathcal{S}\!\mathit{paces}}
\def\Quotfunctor{\mathrm{Quot}}
\def\Hilbfunctor{\mathrm{Hilb}}
\def\Curvesstack{\mathcal{C}\!\mathit{urves}}
\def\Polarizedstack{\mathcal{P}\!\mathit{olarized}}
\def\Complexesstack{\mathcal{C}\!\mathit{omplexes}}
% \Pic is the operator that assigns to X its picard group, usage \Pic(X)
% \Picardstack_{X/B} denotes the Picard stack of X over B
% \Picardfunctor_{X/B} denotes the Picard functor of X over B
\def\Pic{\mathop{\mathrm{Pic}}\nolimits}
\def\Picardstack{\mathcal{P}\!\mathit{ic}}
\def\Picardfunctor{\mathrm{Pic}}
\def\Deformationcategory{\mathcal{D}\!\mathit{ef}}
\usepackage{cleveref}
\crefname{originallemma}{original lemma}{original lemmas}
\Crefname{originallemma}{Original lemma}{Original lemmas}
\makeatletter
\newlabel{coherent-lemma-flat-base-change-cohomology}{{02KH}{source}{}{sourcepointer-1}{}}
\newlabel{coherent-lemma-flat-base-change-cohomology@cref}{{[originallemma][0][]02KH}{[1][0][]source}{}{}{}}
\makeatother
\crefname{originallemma}{original lemma}{original lemmas}
\Crefname{originallemma}{Original lemma}{Original lemmas}
\makeatletter
\newlabel{lemma-affine}{{0245}{source}{}{sourcepointer-2}{}}
\newlabel{lemma-affine@cref}{{[originallemma][0][]0245}{[1][0][]source}{}{}{}}
\makeatother
\crefname{originallemma}{original lemma}{original lemmas}
\Crefname{originallemma}{Original lemma}{Original lemmas}
\makeatletter
\newlabel{lemma-descending-property-affine}{{02L5}{source}{}{sourcepointer-3}{}}
\newlabel{lemma-descending-property-affine@cref}{{[originallemma][0][]02L5}{[1][0][]source}{}{}{}}
\makeatother
\crefname{originallemma}{original lemma}{original lemmas}
\Crefname{originallemma}{Original lemma}{Original lemmas}
\makeatletter
\newlabel{lemma-descending-types-morphisms}{{02W3}{source}{}{sourcepointer-4}{}}
\newlabel{lemma-descending-types-morphisms@cref}{{[originallemma][0][]02W3}{[1][0][]source}{}{}{}}
\makeatother
\crefname{originallemma}{original lemma}{original lemmas}
\Crefname{originallemma}{Original lemma}{Original lemmas}
\makeatletter
\newlabel{morphisms-lemma-base-change-quasi-affine}{{01SO}{source}{}{sourcepointer-5}{}}
\newlabel{morphisms-lemma-base-change-quasi-affine@cref}{{[originallemma][0][]01SO}{[1][0][]source}{}{}{}}
\makeatother
\crefname{originallemma}{original lemma}{original lemmas}
\Crefname{originallemma}{Original lemma}{Original lemmas}
\makeatletter
\newlabel{morphisms-lemma-characterize-quasi-affine}{{01SM}{source}{}{sourcepointer-6}{}}
\newlabel{morphisms-lemma-characterize-quasi-affine@cref}{{[originallemma][0][]01SM}{[1][0][]source}{}{}{}}
\makeatother
\crefname{originallemma}{original lemma}{original lemmas}
\Crefname{originallemma}{Original lemma}{Original lemmas}
\makeatletter
\newlabel{morphisms-lemma-fpqc-quotient-topology}{{02JY}{source}{}{sourcepointer-7}{}}
\newlabel{morphisms-lemma-fpqc-quotient-topology@cref}{{[originallemma][0][]02JY}{[1][0][]source}{}{}{}}
\makeatother
\crefname{originaldefinition}{original definition}{original definitions}
\Crefname{originaldefinition}{Original definition}{Original definitions}
\makeatletter
\newlabel{topologies-definition-fpqc-covering}{{022B}{source}{}{sourcepointer-8}{}}
\newlabel{topologies-definition-fpqc-covering@cref}{{[originaldefinition][0][]022B}{[1][0][]source}{}{}{}}
\makeatother
\title{Descent data for quasi-affine morphisms over an fpqc covering are effective}
\author{The Stacks Project Authors}
\date{stacks-project:descent.tex:local-snapshot-20261012; GNU FDL 1.2 or later}
\begin{document}
\maketitle
\noindent\textit{Editorial note.} All delivered mathematical text is the author's own, copied from the pinned Stacks Project chapter file descent.tex; no statement or proof step was rewritten, reordered or removed.\par\medskip
\noindent\textit{Editorial note.} The counted spans carry navigation stubs for the references leaving them; no bibliography citation occurs in the spans.\par\medskip
\noindent\textit{Editorial note.} The preamble span points at the chapter's own preamble.tex shipped with the frozen source, used verbatim as staged.\par\medskip
\noindent\textit{Allowed standard prerequisites (editorial).}\begin{itemize}
\item Effectivity of descent data for morphisms whose type descends, over an fpqc covering. Same chapter, Lemma lemma-descending-types-morphisms, invoked at proof lines 9465-9468. Reduces the general fpqc covering to the case of a single flat surjective morphism of affines.
\item The characterisation of quasi-affine morphisms by a quasi-compact open immersion into the relative spectrum of the pushforward. Morphisms, Lemmas morphisms-lemma-characterize-quasi-affine and morphisms-lemma-base-change-quasi-affine, invoked at proof lines 9461-9464 and 9475-9480. Produces the open immersion V -> W and the base-change stability used for the reduction.
\item Flat base change for the cohomology of quasi-coherent sheaves. Cohomology of Schemes, Lemma coherent-lemma-flat-base-change-cohomology, invoked at proof lines 9489-9490. Identifies the pullbacks pr\_0\textasciicircum{}* pi\_* O\_V and pr\_1\textasciicircum{}* pi\_* O\_V with the pushforwards of the two base changes of V, so that phi induces phi\textasciicircum{}sharp.
\item Effectivity of descent for affine morphisms, and affineness of the descended scheme. Same chapter, Lemmas lemma-affine and lemma-descending-property-affine, invoked at proof lines 9506-9511. Produces the scheme U' descending W and records that U' is affine.
\item The quotient topology of an fpqc covering. Morphisms, Lemma morphisms-lemma-fpqc-quotient-topology, invoked at proof lines 9528-9529. Shows that the descended subset U of U' is open, which finishes the proof.
\end{itemize}\medskip
\ifdefined\hypertarget\hypertarget{sourcepointer-1}{}\fi
\noindent\textit{Source reference.} Original lemma 02KH: sources/stacks-chow/tags.txt, source lines 3372--3372. Pointer outside the counted span, invoked once for context; the load-bearing argument of the counted statement is the author's proof reproduced in full in the proof span.\par\medskip
\ifdefined\hypertarget\hypertarget{sourcepointer-2}{}\fi
\noindent\textit{Source reference.} Original lemma 0245: sources/stacks-descent/pinned-original-source.tex, source lines 9365--9365. Pointer outside the counted span, invoked once for context; the load-bearing argument of the counted statement is the author's proof reproduced in full in the proof span.\par\medskip
\ifdefined\hypertarget\hypertarget{sourcepointer-3}{}\fi
\noindent\textit{Source reference.} Original lemma 02L5: sources/stacks-descent/pinned-original-source.tex, source lines 6143--6143. Pointer outside the counted span, invoked once for context; the load-bearing argument of the counted statement is the author's proof reproduced in full in the proof span.\par\medskip
\ifdefined\hypertarget\hypertarget{sourcepointer-4}{}\fi
\noindent\textit{Source reference.} Original lemma 02W3: sources/stacks-descent/pinned-original-source.tex, source lines 9245--9245. Pointer outside the counted span, invoked once for context; the load-bearing argument of the counted statement is the author's proof reproduced in full in the proof span.\par\medskip
\ifdefined\hypertarget\hypertarget{sourcepointer-5}{}\fi
\noindent\textit{Source reference.} Original lemma 01SO: sources/stacks-chow/tags.txt, source lines 2345--2345. Pointer outside the counted span, invoked once for context; the load-bearing argument of the counted statement is the author's proof reproduced in full in the proof span.\par\medskip
\ifdefined\hypertarget\hypertarget{sourcepointer-6}{}\fi
\noindent\textit{Source reference.} Original lemma 01SM: sources/stacks-chow/tags.txt, source lines 2343--2343. Pointer outside the counted span, invoked once for context; the load-bearing argument of the counted statement is the author's proof reproduced in full in the proof span.\par\medskip
\ifdefined\hypertarget\hypertarget{sourcepointer-7}{}\fi
\noindent\textit{Source reference.} Original lemma 02JY: sources/stacks-chow/tags.txt, source lines 3353--3353. Pointer outside the counted span, invoked once for context; the load-bearing argument of the counted statement is the author's proof reproduced in full in the proof span.\par\medskip
\ifdefined\hypertarget\hypertarget{sourcepointer-8}{}\fi
\noindent\textit{Source reference.} Original definition 022B: sources/stacks-chow/tags.txt, source lines 2700--2700. Pointer outside the counted span, invoked once for context; the load-bearing argument of the counted statement is the author's proof reproduced in full in the proof span.\par\medskip
\begin{lemma}
\label{lemma-quasi-affine}
Let $S$ be a scheme.
Let $\{X_i \to S\}_{i\in I}$ be an fpqc covering, see
Topologies, Definition \ref{topologies-definition-fpqc-covering}.
Let $(V_i/X_i, \varphi_{ij})$ be a descent datum
relative to $\{X_i \to S\}$. If each morphism
$V_i \to X_i$ is quasi-affine, then the descent datum is
effective.
\end{lemma}

\begin{proof}
Being quasi-affine is a property of morphisms of schemes
which is preserved under any base change, see
Morphisms, Lemmas \ref{morphisms-lemma-characterize-quasi-affine} and
\ref{morphisms-lemma-base-change-quasi-affine}.
Hence Lemma \ref{lemma-descending-types-morphisms} applies
and it suffices to prove the statement of the lemma
in case the fpqc-covering is given by a single
$\{X \to S\}$ flat surjective morphism of affines.
Say $X = \Spec(A)$ and $S = \Spec(R)$ so
that $R \to A$ is a faithfully flat ring map.
Let $(V, \varphi)$ be a descent datum relative to $X$ over $S$
and assume that $\pi : V \to X$ is quasi-affine.

\medskip\noindent
According to Morphisms, Lemma \ref{morphisms-lemma-characterize-quasi-affine}
this means that
$$
V \longrightarrow \underline{\Spec}_X(\pi_*\mathcal{O}_V) = W
$$
is a quasi-compact open immersion of schemes over $X$.
The projections $\text{pr}_i : X \times_S X \to X$ are flat
and hence we have
$$
\text{pr}_0^*\pi_*\mathcal{O}_V =
(\pi \times \text{id}_X)_*\mathcal{O}_{V \times_S X}, \quad
\text{pr}_1^*\pi_*\mathcal{O}_V =
(\text{id}_X \times \pi)_*\mathcal{O}_{X \times_S V}
$$
by flat base change
(Cohomology of Schemes, Lemma \ref{coherent-lemma-flat-base-change-cohomology}).
Thus the isomorphism $\varphi : V \times_S X \to X \times_S V$ (which
is an isomorphism over $X \times_S X$) induces an isomorphism
of quasi-coherent sheaves of algebras
$$
\varphi^\sharp :
\text{pr}_0^*\pi_*\mathcal{O}_V
\longrightarrow
\text{pr}_1^*\pi_*\mathcal{O}_V
$$
on $X \times_S X$.
The cocycle condition for $\varphi$ implies the cocycle condition
for $\varphi^\sharp$. Another way to say this is that it produces
a descent datum $\varphi'$ on the affine scheme $W$ relative to
$X$ over $S$, which moreover has the property that the morphism
$V \to W$ is a morphism of descent data.
Hence by Lemma \ref{lemma-affine}
(or by effectivity of descent for quasi-coherent
algebras) we obtain a scheme $U' \to S$ with an isomorphism
$(W, \varphi') \cong (X \times_S U', can)$ of descent data.
We note in passing that $U'$ is affine by
Lemma \ref{lemma-descending-property-affine}.

\medskip\noindent
And now we can think of $V$ as a (quasi-compact)
open $V \subset X \times_S U'$ with the property that
it is stable under the descent datum
$$
can : X \times_S U' \times_S X \to X \times_S X \times_S U',
(x_0, u', x_1) \mapsto (x_0, x_1, u').
$$
In other words $(x_0, u') \in V \Rightarrow (x_1, u') \in V$
for any $x_0, x_1, u'$ mapping to the same point of $S$.
Because $X \to S$ is surjective we immediately find that
$V$ is the inverse image of a subset $U \subset U'$ under
the morphism $X \times_S U' \to U'$.
Because $X \to S$ is quasi-compact, flat and surjective
also $X \times_S U' \to U'$ is quasi-compact flat and surjective.
Hence by Morphisms, Lemma \ref{morphisms-lemma-fpqc-quotient-topology}
this subset $U \subset U'$ is open and we win.
\end{proof}

\end{document}
